% !TEX program = xelatex
%==============================================================================
% Research Methods - Group assignment - pitch
% Groep 42 - 2025-2026
%==============================================================================

\documentclass[aspectratio=169]{beamer}

% You can increase the font size a bit with the option `14pt'

%==============================================================================
% Preamble
%==============================================================================

%---------- Layout ------------------------------------------------------------

\usetheme{hogent}

% Choose the background color theme, hgwhite (white background) or hgblack
% (black background)
\usecolortheme{hgwhite}

%---------- Packages ----------------------------------------------------------
% You can add additional packages if you need them.

\usepackage[dutch]{babel}      % Dutch language support
\usepackage{booktabs}          % Better looking tables
\usepackage{multirow,multicol} % Merge table cells
\usepackage{eurosym}           % Euro symbool

%---------- Metadata ----------------------------------------------------------

\title{Pitch onderzoeksvoorstellen.}
\subtitle{Research Methods, 25-26, groep 83}
\author{Nicolas Robinet \and Tatjana Schouteeten}
\date{\today}

%==============================================================================
% Presentation content
%==============================================================================

\begin{document}

{
    % HOGENT-huisstijl: afbeelding uitgesneden in de vorm van een letter.
    \setbeamertemplate{background}[imgletter]
        {blurred-background-cellphone-coffee-842554.jpg}{R}

    \begin{frame}
        \maketitle
    \end{frame}
}

\begin{frame}
    \frametitle{Inhoud.}

    \tableofcontents
\end{frame}

%==============================================================================
\section{Webfiltering op het schoolnetwerk.}
%==============================================================================

\begin{frame}
    \frametitle{Waarom dit onderwerp telt.}

    \begin{columns}[t]
        \column{.33\textwidth}
        \textbf{Maatschappelijk}

        \medskip
        Een school heeft een zorgplicht om minderjarigen op haar netwerk
        tegen schadelijke en leeftijdsongeschikte content te beschermen.

        \column{.33\textwidth}
        \textbf{Binnen de specialisatie}

        \medskip
        DNS-filtering, firewall-categorisatie, TLS-inspectie en
        netwerkbeheer vormen de kern van System \& Network Engineering.

        \column{.33\textwidth}
        \textbf{Actueel}

        \medskip
        Versleutelde DNS-protocollen maken klassieke DNS-filtering steeds
        eenvoudiger omzeilbaar.
    \end{columns}
\end{frame}

\subsection{Probleemstelling.}

\begin{frame}
    \frametitle{Probleemstelling.}

    College Ten Doorn wil contentfiltering op haar schoolnetwerk
    \textbf{sluitend afdwingen} voor haar minderjarige leerlingen.

    \bigskip

    De huidige, beperkte filtering is niet centraal geëvalueerd.
    Klassieke DNS-filtering kan bovendien worden omzeild via
    DNS-over-HTTPS, alternatieve resolvers en VPN-verbindingen.

    \bigskip
    \bigskip

    \textcolor{hggrey}{\textbf{Onderbouwd door bronnen.}} Hoffman \&
    McManus (2018) en Pearce et al.\ (2017) beschrijven de beperkingen en
    omzeilingsroutes van DNS-filtering.
\end{frame}

\subsection{Onderzoeksvragen.}

\begin{frame}
    \frametitle{Hoofdonderzoeksvraag.}

    \begin{quote}
        {\large Welke webfilteroplossing is het meest geschikt om het
        contentfilterbeleid van College Ten Doorn op haar schoolnetwerk af
        te dwingen, rekening houdend met haar infrastructuur, budget en
        beheerscapaciteit?}
    \end{quote}

    \bigskip

\end{frame}

\begin{frame}
    \frametitle{Deelvragen.}

    \begin{columns}[t]
        \column{.5\textwidth}
        \textbf{Probleemdomein}
        \begin{itemize}
            \item Hoe zien de huidige netwerkinfrastructuur en het
                  filterbeleid eruit, en waar schieten die tekort?
            \item Aan welke functionele en niet-functionele requirements
                  moet een oplossing voldoen?
            \item Welke webfiltertechnieken bestaan er en wat zijn hun
                  beperkingen en omzeilingsmogelijkheden?
        \end{itemize}

        \column{.5\textwidth}
        \textbf{Oplossingsdomein}
        \begin{itemize}
            \item Welke oplossingen komen op de longlist en welke voldoen
                  het best aan de requirements (shortlist)?
            \item Hoe presteert de beste oplossing op blokkeringsratio,
                  vals-positieven, latency en doorvoer?
            \item In welke mate weerstaat ze DNS-over-HTTPS, alternatieve
                  DNS-servers en VPN-verbindingen?
        \end{itemize}
    \end{columns}
\end{frame}

\subsection{Methodologie.}

\begin{frame}
    \frametitle{Methodologie: probleemdomein.}

    \footnotesize
    \begin{table}
        \renewcommand{\arraystretch}{1.4}
        \begin{tabular}{@{}lp{2.7cm}p{4cm}p{3.3cm}@{}}
            \toprule
            & \textbf{Methode} & \textbf{Databronnen / input} &
            \textbf{Output} \\
            \midrule
            \textbf{DV1} & Casusanalyse &
            Interview IT-dienst en analyse van het netwerkplan &
            Infrastructuur, filterbeleid en pijnpunten \\
            \midrule
            \textbf{DV2} & Interviews + MoSCoW &
            IT-dienst en directie &
            Geprioriteerde functionele en niet-functionele requirements \\
            \midrule
            \textbf{DV3} & Literatuurstudie &
            Documentatie over DNS-filtering, proxy-/TLS-inspectie en
            firewall-categorisatie &
            Technieken, beperkingen en omzeilingsroutes \\
            \bottomrule
        \end{tabular}
    \end{table}
\end{frame}

\begin{frame}
    \frametitle{Methodologie: oplossingsdomein.}

    \footnotesize
    \begin{table}
        \renewcommand{\arraystretch}{1.4}
        \begin{tabular}{@{}lp{2.7cm}p{4cm}p{3.3cm}@{}}
            \toprule
            & \textbf{Methode} & \textbf{Databronnen / input} &
            \textbf{Output} \\
            \midrule
            \textbf{DV4} & Vergelijking &
            Kandidaat-oplossingen &
            Longlist en shortlist \\
            \midrule
            \textbf{DV5} & Proof of concept &
            VM-testomgeving en testcorpus &
            Blokkering, vals-positieven, latency en doorvoer \\
            \midrule
            \textbf{DV6} & Omzeilingstest &
            DNS-over-HTTPS, alternatieve DNS en VPN &
            Geslaagde bypasses \\
            \bottomrule
        \end{tabular}
    \end{table}

    \bigskip

    \textcolor{hggrey}{\textbf{Eigen technische bijdrage.}}
    Reproduceerbare testopstelling; data en scripts.
\end{frame}

%==============================================================================
\section{Digitaal toezicht in de jeugdzorg.}
%==============================================================================

\begin{frame}
    \frametitle{Waarom dit onderwerp telt.}

    \begin{columns}[t]
        \column{.33\textwidth}
        \textbf{Maatschappelijk}

        \medskip
        Minderjarige bewoners van een zorginstelling hebben een
        beschermend kader nodig voor internettoegang en gamen.

        \column{.33\textwidth}
        \textbf{Interprofessioneel}

        \medskip
        Het thema verbindt netwerkbeheer met zorgcoördinatie,
        pedagogisch beleid en het welzijn van een kwetsbare doelgroep.

        \column{.33\textwidth}
        \textbf{Actueel}

        \medskip
        Grooming, cyberpesten, schadelijke content en gokmechanismen in
        games blijven reële online risico's.
    \end{columns}
\end{frame}

\subsection{Probleemstelling.}

\begin{frame}
    \frametitle{Probleemstelling.}

    Zorginstelling Y is een jeugdpsychiatrische voorziening in
    Oost-Vlaanderen. Zij moet afspraken over schermtijd, leeftijdsgeschikte
    content en veilige communicatie \textbf{centraal op het netwerk
    afdwingen}.

    \bigskip

    PEGI-classificaties informeren, maar dwingen niets af. Ingebouwde
    parental controls werken per toestel en platform en zijn voor een
    instelling met wisselende toestellen onbeheersbaar.

    \bigskip
    \bigskip

    \textcolor{hggrey}{\textbf{Onderbouwd door bronnen.}} Smahel et al.\
    (2020), Mediawijs (2024), Carville \& Chen (2024) en UNICEF (2024)
    documenteren de online risico's voor minderjarigen.
\end{frame}

\subsection{Onderzoeksvragen.}

\begin{frame}
    \frametitle{Hoofdonderzoeksvraag.}

    \begin{quote}
        {\large Welke netwerkgebaseerde oplossing voor digitaal toezicht
        is het meest geschikt om het internet- en gamebeleid van
        zorginstelling Y voor haar minderjarige bewoners af te dwingen?}
    \end{quote}

    \bigskip

    
\end{frame}

\begin{frame}
    \frametitle{Deelvragen.}

    \begin{columns}[t]
        \column{.5\textwidth}
        \textbf{Probleemdomein}
        \begin{itemize}
            \item Welke online risico's binnen games en platformen zijn
                  het meest relevant voor de bewoners?
            \item Welk internet- en gamebeleid wil de instelling
                  afdwingen, en welke requirements volgen daaruit?
            \item Waarom volstaan PEGI en ingebouwde parental controls
                  niet op het niveau van een instelling?
        \end{itemize}

        \column{.5\textwidth}
        \textbf{Oplossingsdomein}
        \begin{itemize}
            \item Welke toezichtoplossingen komen in aanmerking en welke
                  voldoen het best aan de requirements?
            \item Hoe accuraat identificeert en reguleert de beste
                  oplossing game- en chatverkeer en tijdschema's?
            \item Welke impact heeft de oplossing op toegelaten verkeer,
                  en in welke mate weerstaat ze versleutelde DNS of VPN?
        \end{itemize}
    \end{columns}
\end{frame}

\subsection{Methodologie.}

\begin{frame}
    \frametitle{Methodologie: probleemdomein.}

    \footnotesize
    \begin{table}
        \renewcommand{\arraystretch}{1.4}
        \begin{tabular}{@{}lp{2.7cm}p{4cm}p{3.3cm}@{}}
            \toprule
            & \textbf{Methode} & \textbf{Databronnen / input} &
            \textbf{Output} \\
            \midrule
            \textbf{DV1} & Casus- en risicoanalyse &
            Literatuurstudie; interview met IT-verantwoordelijke en
            zorgcoördinatie &
            Risicoprofiel van bewonersgroep en huidige netwerkopzet \\
            \midrule
            \textbf{DV2} & Interviews + MoSCoW &
            IT-verantwoordelijke en zorgcoördinatie &
            Geprioriteerde functionele en niet-functionele requirements \\
            \midrule
            \textbf{DV3} & Literatuurstudie &
            PEGI en ingebouwde parental controls &
            Lokale controls tekort \\
            \bottomrule
        \end{tabular}
    \end{table}
\end{frame}

\begin{frame}
    \frametitle{Methodologie: oplossingsdomein.}

    \footnotesize
    \begin{table}
        \renewcommand{\arraystretch}{1.4}
        \begin{tabular}{@{}lp{2.7cm}p{4cm}p{3.3cm}@{}}
            \toprule
            & \textbf{Methode} & \textbf{Databronnen / input} &
            \textbf{Output} \\
            \midrule
            \textbf{DV4} & Vergelijking &
            Kandidaat-oplossingen &
            Longlist en shortlist \\
            \midrule
            \textbf{DV5} & Proof of concept &
            Games, chatplatformen en schooltoepassingen &
            Herkenning, vals-positieven en tijdschema's \\
            \midrule
            \textbf{DV6} & Prestatie- en omzeilingstest &
            Toegelaten verkeer, DNS-over-HTTPS en VPN &
            Latency, doorvoer en bypasses \\
            \bottomrule
        \end{tabular}
    \end{table}

    \bigskip

    \textcolor{hggrey}{\textbf{Eigen technische bijdrage.}}
    Test van verkeer, tijdschema's en omzeiling.
\end{frame}

%==============================================================================
\section{Bronnen.}
%==============================================================================

\begin{frame}
    \frametitle{Bronnen: webfiltering (APA).}

    \small
    \begin{itemize}
        \item Hoffman, P., \& McManus, P. (2018). \textit{DNS Queries
              over HTTPS (DoH)} (RFC 8484). IETF.
        \item Pearce, P., et al.\ (2017). \textit{Global Measurement of
              DNS Manipulation}. USENIX Security Symposium.
        \item Pi-hole. (2026). \textit{Pi-hole documentation}.
              \url{https://docs.pi-hole.net}
        \item Netgate. (2026). \textit{pfSense Documentation}.
              \url{https://docs.netgate.com/pfsense}
        \item Cisco. (2026). \textit{Cisco Umbrella documentation}.
              \url{https://docs.umbrella.com}
    \end{itemize}

    \medskip
    \textcolor{hggrey}{\footnotesize Bronnen voor de technische
    vergelijking, de omzeilingsscenario's en kandidaat-oplossingen.}
\end{frame}

\begin{frame}
    \frametitle{Bronnen: digitaal toezicht (APA).}

    \small
    \begin{itemize}
        \item Smahel, D., et al.\ (2020). \textit{EU Kids Online 2020:
              Survey results from 19 countries}. EU Kids Online.
        \item Mediawijs \& imec. (2024). \textit{Apestaartjaren}.
              \url{https://www.apestaartjaren.be}
        \item Carville, O., \& Chen, C. D. (2024). \textit{Roblox's
              Predator Problem}. Bloomberg Businessweek, 4822, 58--67.
        \item Zendle, D., \& Cairns, P. (2018). \textit{Video game loot
              boxes are linked to problem gambling}. PLOS ONE.
        \item UNICEF. (2024). \textit{Keeping children safe online};
              European Commission. (2022). \textit{A better internet for
              kids (BIK+)}.
    \end{itemize}
\end{frame}

%==============================================================================

{
\setbeamertemplate{background}[imgletter]
    {blurred-background-cellphone-coffee-842554.jpg}{B}

\begin{frame}

    {\huge \textbf{Bedankt. Vragen?}}

    \bigskip

    \textbf{Groep 83} \\
    Nicolas Robinet, Tatjana Schouteeten

\end{frame}
}

\end{document}
